% Focal incorporation of prior dependencies; does not replace the common kernel.
% Detailed provenance: technical/CIERRE_DEPENDENCIAS_ANTERIORES_VI.md.
% Test values are published after their operators; they are not inputs.
\section{Generated coordinates, action, and electronic normalisation}
\label{vi:dep:seccion}\label{sec:k-registro}

The construction of the collections uses results from the preceding stages:
the regional characters, the dodecaphase record, the action line, and the
electronic section. We assemble here the operations connecting them and the
proofs needed to compose them. Their order is causal: the regions and their
records are produced on APP through TRIT and TPK; real evaluations
are performed afterwards; the choice of units enters when dimensional
magnitudes are realised.

Initial conditions and regional selection were developed in
the section preceding the arithmetic of histories. The domain of the incidence
ledger, its counts, and the reconstruction of the record are defined
explicitly below. Selecting a specific collection of
visits and traces is an additional operation beyond evaluating a ledger:
it requires specifying its multiplicities, boundary transport, and
coverage of the admitted traces. The reconstruction formulas
proved here allow that record to be evaluated and checked; their
invertibility does not by itself select the incidence collection.

\subsection{Regional readers and compatibility of evaluation}
\label{vi:dep:regiones}

Closure, propagation, and self-scaling are three regions of the same
generation. The closure region has five emissions and multiplicity
\(432+4\cdot144=1008\). In the oriented gauge of the catalogue, its emissions
are
\[
\begin{array}{c|r}
(501,614,498,169,272,272)&432\\
(810,923,870,169,272,272)&144\\
(810,983,810,169,272,272)&144\\
(870,923,810,169,272,272)&144\\
(870,923,810,418,674,521)&144.
\end{array}
\]
The common ternary word is \(010211\); the representatives of propagation
and self-scaling are \(201101\) and \(121200\). The census and orientation
precede the evaluations that follow.

At the five positions of the closure reader, the uniform projection
\(P_5=\mathbf1\mathbf1^{\mathsf T}/5\), together with
\(P_5\lambda=\lambda\) and \(\mathbf1^{\mathsf T}\lambda=1\), determines
\(\lambda=\mathbf1/5\). This normalisation of the reader positions
differs from the seed distribution given in the table.
In the elliptic regime \(J^2=-I\), multiplication gives
\[
(I+aJ)(I+bJ)=(1-ab)I+(a+b)J.
\]
Whenever \(1-ab\ne0\), the composite slope is
\(a\oplus b=(a+b)/(1-ab)\). Two doublings of \(1/5\) give
\(5/12\) and \(120/119\). The positive compensator \(1/q\), \(q>1\),
that returns the slope to one satisfies
\[
\frac{120/119-1/q}{1+(120/119)/q}=1,
\qquad q=239.
\]
Thus the closure character of this reader is
\begin{equation}
 \Pi_{\rm cl}=4\left(4\arctan\frac15-\arctan\frac1{239}\right).
 \label{vi:dep:lector-clausura}
\end{equation}
The composition rule, its compensation, and the choice of angular branch
fix this value; the conventional angular identity recognises it as \(\pi\).

The propagation character is realised in a commutative formal collection
\(P_t\), normalised by \(P_0=1\), \(P'_0=1\), and
\(P_{t+s}=P_tP_s\). If \(P_t=\sum_{n\ge0}a_nt^n\), comparison
of the coefficient of \(s^nt\) gives
\[
 (n+1)a_{n+1}=a_n,\qquad a_0=1.
\]
Hence \(a_n=1/n!\); representation at one unit of the parameter is
\(\Pi_{\rm pr}=\sum_{n\ge0}1/n!\), recognised as \(e\).
The unit and the sign of the generator belong to the reader: the semigroup
law alone would admit other generators.

For self-scaling, TRIT selects the modes
\((\Sigma,\Pi,\Pi)\). Their three cyclic transitions are
\(\Sigma\to\Pi\), \(\Pi\to\Pi\), and \(\Pi\to\Sigma\). Their incidence
matrix, in areal--volumetric order, is
\begin{equation}
 F=\begin{pmatrix}0&1\\1&1\end{pmatrix}.
 \label{vi:dep:incidencia-hojas}
\end{equation}
The positive eigenvector normalised as \((1,x)^{\mathsf T}\)
satisfies \(x=\lambda\), \(1+x=\lambda x\); hence
\(x^2-x-1=0\). Its positive root is the self-scaling coordinate
\(\varphi\). This equation recognises the character of the incidence already
produced by the calendar.

\begin{proposition}[Regional evaluation through compatible intervals]
\label{vi:dep:prefijos}
The three preceding readers admit rational intervals of arbitrarily
small width. Their finite decimal prefixes are determined by
refinement and are compatible under truncation.
\end{proposition}
\begin{proof}
For \(0<x<1\), two consecutive alternating partial sums of
\(\sum_{j\ge0}(-1)^jx^{2j+1}/(2j+1)\) enclose \(\arctan x\).
The first omitted term bounds the error. This provides rational
intervals for \eqref{vi:dep:lector-clausura}.
If \(E_n=\sum_{j=0}^n1/j!\), then, for \(n\ge1\),
\[
 E_n<\Pi_{\rm pr}<E_n+
 \frac{n+2}{(n+1)(n+1)!}.
\]
The bound follows by bounding successive ratios of the tail by
\(1/(n+2)\). Finally, \(F^n\) has entries
\(F_{n-1},F_n,F_{n+1}\), where \(F_0=0,F_1=1\) and
\(F_{n+1}=F_n+F_{n-1}\). The identity
\(F_{n+1}^2-F_nF_{n+2}=(-1)^n\) shows that two consecutive ratios
enclose \(\varphi\) with width \(1/(F_nF_{n+1})\).
Successive intervals may be intersected to preserve
nestedness. After conventional recognition, the irrationality of
\(\pi,e,\varphi\) ensures that no coordinate lies on a rational
decimal boundary. At each finite depth, a sufficiently
narrow interval lies in a unique decimal cell. Its index
determines the prefix, and nestedness proves compatibility.
\end{proof}

\subsection{Signed dodecaphase chart of the full TPK state}
\label{vi:dep:estado-pleno}

The evolution producing the record is formulated on the complete state of the
Topological Phase Kernel. Denote by
\(\mathsf{TPKFullState}\) the state space jointly preserving
APP position, TRIT regime, route, sheet, orientation, carry, boundary,
memory, incidence ledger, and terminal coordinates. In a canonical chart,
\[
 X=(X_{\rm APP},X_{\rm TRIT},X_{\rm ruta},X_{\rm mem},
       C_{\rm term},U_{12}^{\rm sgn},X_{\partial},\ldots).
\]
The dynamics of one hundred and eight steps acts through the composition
\begin{equation}
 X_{\rm term}
 =\bigl(\mathcal U_{108}\circ\cdots\circ\mathcal U_1\bigr)(X_{\rm seed}),
 \qquad
 \mathcal U_t=\mathsf{Upd}_t\circ\mathsf{Tra}_t\circ\mathsf{Sel}_t,
 \label{vi:dep:estado-terminal-108}
\end{equation}
where \(\mathsf{Sel}_t\) applies the tritic phase selector,
\(\mathsf{Tra}_t\) transports sheet, route, and orientation, and
\(\mathsf{Upd}_t\) incorporates the carry and memory produced at
instant \(t\). Each factor therefore acts on the full state, and the
composition preserves the genealogy of the twelve nonadic windows.

The incidence ledger of the following section is the coordinate
\(\mathcal R_{12}(X_{\rm term})\). Its evaluation by the
\(90/120\) channels defines
\[
 \mathcal E_{108}^{90,120}\circ\mathcal R_{12}:
 \mathsf{TPKFullState}\longrightarrow\mathbb Z^{25},
 \qquad
 X\longmapsto C(X)=(b^{90}(X),b^{120}(X),Q(X)).
\]
The signed dodecaphase projection is
\[
 \mathcal S_{12}:=\operatorname{pr}_{U_{12}^{\rm sgn}}:
 \mathsf{TPKFullState}\longrightarrow\mathbb Z^{12}.
\]
On the compatible sublattice produced by the dynamics, the two charts
satisfy the operator identity
\begin{equation}
 \mathcal S_{12}
 =\Pi_H\circ\mathcal E_{108}^{90,120}\circ\mathcal R_{12}
 =\Pi_H\circ\operatorname{pr}_{C}.
 \label{vi:dep:identidad-productor-e108}
\end{equation}
Thus the dodecaphase signature is obtained by projection and evaluation of the same
state that preserves the ledger; its domain includes the orientation and memory
entering each count. The raw projection
\(\rho_{108}:\mathsf{TPKFullState}\to\mathsf{CT108RawProjection}\),
by contrast, denotes the chart that forgets those coordinates after
production. This separation of types fixes the causal order between generation,
signed representation, and reduced reading.

\subsection{From temporal incidences to the dodecaphase record}
\label{vi:dep:registro}\label{subsec:k-registro-integral}

\begin{definition}[Domain of the oriented incidence ledger]
\label{vi:dep:libro-dominio}
A ledger \(\mathscr L\) contains a nonempty finite collection of routes,
their identified visits, and a finite collection of traces. Each route
preserves one visit per instant during the first \(108\) instants,
and the later visits used by its traces. A visit preserves
\[
 v=(\gamma,t,i,j,s,\tau,w,\varepsilon_\partial),
\]
where \(\gamma\) identifies the route, \(t\) its instant, \(1\le i,j\le9\)
its APP position, \(s\in\{\Sigma,\Pi,0\}\) its sheet,
\(\tau\in\{-1,0,1\}\) its tritic orientation, and \(w\) its positive
integer multiplicity. The incidence chart evaluates
\[
 n_v=
 \begin{cases}
 i+j,&s=\Sigma,\\
 ij,&s=\Pi,\\
 9,&s=0.
 \end{cases}
 \qquad
 r_v=1+((n_v-1)\bmod9),\quad q_v=(n_v-r_v)/9.
\]
The reading \(n_v=9\) fixes the threshold convention of this chart.
For \(r_v=9\), the label \(\varepsilon_\partial\) distinguishes
corona from interior; it is not obtained from the remainder alone.

A trace \(\theta=(v_1,\ldots,v_N)\) traverses consecutive instants of
the same route. It preserves a multiplicity \(u(\theta)\in\mathbb N_{>0}\)
and a reduced length unit
\(\lambda(\theta)\in\mathbb N_{>0}\), with
\[
 N\lambda(\theta)=120(1+3j(\theta)),\qquad j(\theta)\in\mathbb N_0.
\]
Its channel is
\(\sigma(\theta)=\operatorname{sgn}(\sum_{a=1}^N\tau(v_a))\).
The collection counted in the \(120\) channels comprises traces of
nonzero sign. Reduced length and number of instants are
related through \(\lambda\); they are not identified with one another.
\end{definition}

The definition specifies the evaluation domain, including the
data that must be transmitted to the reader. Its multiplicities and trace
collection belong to the dynamic selection of the ledger. An enumeration
of the domain or a format check does not replace that selection.

Generation preserves a ledger \(\mathscr L\) of visits and traces.
The temporal chart divides it into windows
\(E_m=\{9m-8,\ldots,9m\}\), \(1\le m\le12\),
equivalent to positions \(m=\beta+1\) of the record
\(\mathcal R_{12}\). Each window retains its subroute,
orientation, boundary increment, and local incidences.
Each visit records APP position, arithmetic sheet, instant, multiplicity,
and boundary orientation. Each evaluation \(n_v\) preserves the
division \(n_v=r_v+9q_v\), with \(1\le r_v\le9\). In window
\(m\), the incidence reader expresses
\begin{align}
 A_m&=\sum_{v\in\mathscr L_m}w(v)
      \mathbf1_{\{1,3,5,7\}}(r_v),\notag\\
 C_m&=\sum_{j\ge0}
       \bigl(\nu^-_{120,m}(j)-\nu^+_{120,m}(j)\bigr),\notag\\
 V_m&=\sum_{v\in\mathscr L_m}w(v)
      \mathbf1_{\{9\}}(r_v)\,\varepsilon_{\partial}(v).
 \label{vi:dep:incidencias}
\end{align}
Here the traces in each window have finite support,
\(\varepsilon_{\partial}=+1\) on the corona and \(-1\) in the interior,
and the labels \(\pm,120\) belong to the preserved transport.
The multiplicity \(w\) and the traces \(\nu\) are those obtained from that
transport; they are not reconstructed by selecting a value of \(\alpha\).

In terms of the identified traces, the second sum has the finite
expression
\[
 C_m=-\sum_{\theta:\,t(v_1)\in E_m}u(\theta)\sigma(\theta).
\]
Indeed, grouping by \(j(\theta)\) and by sign produces
\(\nu^-_{120,m}(j)-\nu^+_{120,m}(j)\). This equality specifies
the operation actually applied to each record; it keeps
channel, multiplicity, and length separate.

\begin{proposition}[Updating and concatenation of the ledger]
\label{vi:dep:libro-actualizacion}
Once the collection of visits and traces admitted by transport is fixed, its
counts are updated through integer increments. When a
visit \(v\) in window \(m\) is added, the increment of \((A_m,C_m,V_m)\) is
\[
 \left(w(v)\mathbf1_{\{1,3,5,7\}}(r_v),\,0,\,
 w(v)\mathbf1_{\{9\}}(r_v)\varepsilon_\partial(v)\right).
\]
When a trace \(\theta\) whose origin belongs to \(E_m\) is completed,
the increment is \((0,-u(\theta)\sigma(\theta),0)\).
For two consecutive segments, the count of their concatenation is
the sum of the counts of the visits and traces contained in each
segment, plus that of the traces crossing their junction.
\end{proposition}
\begin{proof}
Each sum in \eqref{vi:dep:incidencias} is indexed by visit
or trace identifiers, so adding an identifier adds
exactly the summand displayed. For concatenation, visits
are distributed between the two segments. Traces are partitioned into
those contained in the first, those contained in the second, and those
crossing the junction. Additivity of integer sums proves
the formula. Preserving the identifiers and pending
prolongations prevents a trace from being counted twice or one that has
not yet ended from being lost. Decimal reduction is applied after
counting; the quotients of that reduction remain in memory.
\end{proof}

Decimal reduction of these three counts produces
\begin{equation}
 z_m=100[A_m]_{10}+10[C_m]_{10}+[V_m]_{10},
 \qquad 0\le z_m<1000.
 \label{vi:dep:bloque-incidencial}
\end{equation}
Let \((D_3z)_i=z_i-z_{i+3}\),
\((D_4z)_i=z_i-z_{i+4}\), with indices in \(\ZZ/12\ZZ\), and
\(Q(z)=\sum_i z_i\). The reader preserves
\(\mathcal D(z)=(D_3z,D_4z,Q(z))\). Decimal reduction forgets
multiples of ten in each count, whereas the original ledger preserves
the visits and transports that generated them.
The reading of even-code events on that same history,
its decomposition \(12=1+5+6\), and its integral inversion are
developed in Proposition \ref{vi:prop:eventos}. That event
vector does not replace the present incidence record:
the domains and preserved coordinates remain distinct.

\begin{proposition}[Incidence conjugation and decimal reduction]
\label{vi:dep:conjugacion-incidencial}
Let \(\varrho(m)=13-m\). Consider a conjugation of the ledger
that maps the visits and traces of window \(m\) to those of
\(\varrho(m)\) through bijections preserving their multiplicities.
Suppose it preserves the arithmetic class and boundary label
of the visits, and interchanges the positive and negative channels of the
traces. Then the counts in \eqref{vi:dep:incidencias} satisfy
\begin{equation}
 (A_m^\dagger,C_m^\dagger,V_m^\dagger)
 =(A_{\varrho(m)},-C_{\varrho(m)},V_{\varrho(m)}).
 \label{vi:dep:conjugacion-recuentos}
\end{equation}
Writing \(u_m=[C_m]_{10}\), their decimal blocks satisfy
\begin{equation}
 z_m^\dagger=z_{\varrho(m)}
       +100\mathbf1_{\{u_{\varrho(m)}\ne0\}}-20u_{\varrho(m)}.
 \label{vi:dep:conjugacion-bloque}
\end{equation}
\end{proposition}
\begin{proof}
The bijections preserve every summand of \(A_m\) and \(V_m\).
Interchanging the channels permutes \(\nu^-\) and \(\nu^+\),
thereby reversing the sign of \(C_m\). This yields
\eqref{vi:dep:conjugacion-recuentos}. If \(C=10q+u\),
\(0\le u<10\), Euclidean division of \(-C\) gives
\[
 [-C]_{10}=10\mathbf1_{\{u\ne0\}}-u,
 \qquad
 \left\lfloor\frac{-C}{10}\right\rfloor
 =-q-\mathbf1_{\{u\ne0\}}.
\]
In \eqref{vi:dep:bloque-incidencial}, only the middle digit
changes. Its difference is
\(10([-C]_{10}-[C]_{10})=100\mathbf1_{\{u\ne0\}}-20u\),
which proves \eqref{vi:dep:conjugacion-bloque}. The second
identity simultaneously preserves the integer quotient of the reduction.
\end{proof}

The proposition specifies the action on the ledger before transporting
its decimal representation. Its application to route reversal
also preserves the convention used to read the endpoints. Let
\(\gamma=(x_0,\ldots,x_n)\) be an enriched route and let \(f\) be an
integer observable evaluated before each advance. The reverse route
\(\gamma^{-1}=(x_n,\ldots,x_0)\) satisfies
\begin{equation}
 \sum_{k=0}^{n-1}f(x_{n-k})
 =\sum_{k=0}^{n-1}f(x_k)+f(x_n)-f(x_0).
 \label{vi:dep:conjugacion-frontera}
\end{equation}
Indeed, the left-hand side sums indices \(1,\ldots,n\);
the first term on the right sums \(0,\ldots,n-1\). Their difference
is exactly the stated boundary term. The identity applies
to each incidence observable with its preserved data. The endpoints
are accounted for in the integer counts before reduction modulo ten;
on a closed route the term vanishes, whereas in an open
window it can change the block. Traces crossing a junction
are preserved through the partition in
Proposition~\ref{vi:dep:libro-actualizacion}.

This connects the counts to route conjugation in
Section~\ref{vi:sec:conjugacion}, with the specific action
on each incidence kept explicit. Negating a channel and reversing
a route are operations on the ledger; their decimal reduction is the
representation just calculated.

\begin{proposition}[Integral image and reconstruction of the record]
\label{vi:dep:imagen-integral}
Let \(b,c\in\ZZ^{12}\), \(q\in\ZZ\), with indices from zero to eleven.
Define
\[
 w_i=c_i-b_{i+1},\qquad P_0=0,\qquad
 P_i=\sum_{j=0}^{i-1}w_j,\qquad T_w=\sum_{i=0}^{11}P_i.
\]
There exists a unique \(k\in\ZZ^{12}\) such that
\(D_3k=b\), \(D_4k=c\), \(Q(k)=q\), if and only if
\begin{equation}
 b_i=w_i+w_{i+1}+w_{i+2},\qquad
 \sum_iw_i=0,\qquad q+T_w\equiv0\pmod{12}.
 \label{vi:dep:condiciones-imagen}
\end{equation}
In that case,
\begin{equation}
 k_i=\frac{q+T_w}{12}-P_i.
 \label{vi:dep:inversa-incidencial}
\end{equation}
\end{proposition}
\begin{proof}
If \(k\) exists, then
\(c_i-b_{i+1}=k_i-k_{i+1}\). The cyclic sum of these differences
is zero, their sum over length three is \(b_i\), and
\(k_i=k_0-P_i\). Summing the twelve coordinates gives
\(q=12k_0-T_w\), proving the necessary conditions and the formula.
Conversely, the conditions make the formula integral and
cyclic. Its differences over length three are \(b_i\), and those over
length four are \(w_i+b_{i+1}=c_i\). Its sum is \(q\).
The difference between two solutions has all consecutive
differences zero and sum zero; it is therefore zero.
\end{proof}

Reconstruction is an operation on the typed image, not a
retrospective selection of its arguments. In particular,
\(\mathcal D(z)=\mathcal D(z')\) implies \(z=z'\).
Since the three-digit expression of an integer in \([0,999]\) is unique,
two triples of counts express the same record if and only if they coincide
componentwise modulo ten.

The signed change of representation uses the block
\begin{equation}
 H_4=\begin{pmatrix}
 1&1&1&1\\1&1&-1&-1\\
 1&-1&1&-1\\1&-1&-1&1
 \end{pmatrix},\qquad H_4^2=4I_4.
 \label{vi:dep:hadamard}
\end{equation}
The operator \(H_{12}\) applies \(H_4\) to the three ordered orbits
\((1,4,7,10)\), \((2,5,8,11)\), \((3,6,9,12)\).
On the sublattice
\(\mathcal L_H=\{U\in\ZZ^{12}:H_{12}U\equiv0\pmod4\}\),
the maps
\begin{equation}
 U=H_{12}K,\qquad K=\tfrac14H_{12}U
 \label{vi:dep:K-reversible}
\end{equation}
are inverses. The congruence expresses exactly the integrality of the
second map, and \(H_{12}^2=4I\) proves both compositions.

The signed record is also calculated directly from the transverse
coordinates before reconstructing \(K\). With indices from zero to eleven,
let, for \(r=0,1,2\),
\[
 S_r=\sum_{j=0}^3c_{r+3j},\quad
 a_0=\frac{q+2S_0+S_1}{3},\quad
 a_1=a_0-S_0,\quad a_2=a_1-S_1,
 \qquad d_{r,j}=b_{r+3j}.
\]
The harmonic map expresses by orbits
\begin{equation}
 \Pi_H^{(r)}(b,c,q)=
 (a_r,d_{r,1}-d_{r,3},d_{r,0}+d_{r,2},d_{r,0}-d_{r,2}).
 \label{vi:dep:PiH}
\end{equation}
To prove the formula on the integral image, let \(k\) be its
preimage. If \(s_r=\sum_jk_{r+3j}\), then
\(S_r=s_r-s_{r+1}\) and \(q=s_0+s_1+s_2\), with orbit indices
modulo three. Solving gives \(s_r=a_r\). For the four coordinates
\((x_0,x_1,x_2,x_3)\) of an orbit,
\(d_j=x_j-x_{j+1}\). Their three combinations in
\eqref{vi:dep:PiH} are respectively
\(x_0+x_1-x_2-x_3\), \(x_0-x_1+x_2-x_3\), and
\(x_0-x_1-x_2+x_3\). Together with \(a_r=\sum_jx_j\), they are
exactly \(H_4(x_0,x_1,x_2,x_3)^{\mathsf T}\).
The effective composition is thus
\(\mathscr L\to(b,c,q)\to\Pi_H(b,c,q)=U\to H_{12}U/4=K\).
The preimage used in this proof verifies the identity;
it is not an argument of the map \(\Pi_H\).

For the terminal state of \eqref{vi:dep:estado-terminal-108}, the coordinate
produced before signed representation is
\begin{align}
 b^{90}={}&(-495,-116,-684,108,601,-90,
              -173,-88,313,560,-397,461),\notag\\
 b^{120}={}&(-425,-281,-481,671,-255,30,
               475,-543,680,251,6,-128),\notag\\
 Q_{\rm TPK}={}&6263.
 \label{vi:dep:canales-terminales}
\end{align}
Substituting these twenty-five fields into \eqref{vi:dep:PiH}, the three
signed orbits are
\begin{align*}
 U_{\rm term}^{(1)}&=(2378,-452,-668,-322),\\
 U_{\rm term}^{(2)}&=(1406,998,-204,-28),\\
 U_{\rm term}^{(3)}&=(2479,-551,-371,-997).
\end{align*}
Interleaving by columns produces
\begin{equation}
 U_{\rm term}=\mathcal S_{12}(X_{\rm term})
 =(2378,1406,2479,-452,998,-551,-668,-204,-371,-322,-28,-997).
 \label{vi:dep:U-terminal}
\end{equation}
The orbital transform \(H_{12}/4\) is defined on this sublattice by
the exact divisibility of every block of four. Consequently, the chain
\begin{equation}
 X_{\rm term}\xrightarrow{\mathcal R_{12}}\mathscr L(X_{\rm term})
 \xrightarrow{\mathcal E_{108}^{90,120}}C_{\rm term}
 \xrightarrow{\Pi_H}U_{\rm term}
 \xrightarrow{H_{12}/4}K
 \label{vi:dep:cadena-terminal-completa}
\end{equation}
is a typed composition from the full state to its periodic seal.
Propositions~\ref{vi:dep:libro-actualizacion} and
\ref{vi:dep:imagen-integral} provide, respectively, the finite update
of the ledger and the integral inverse of the last representation.

The generated record provides the ordered witness
\begin{equation}
 K=(234,543,140,729,659,824,621,058,914,794,146,601).
 \label{vi:dep:K-testigo}
\end{equation}
The leading zero in \(058\) forms part of the block reading.
To check it through Proposition~\ref{vi:dep:imagen-integral}, its differences
and sum are
\[
\begin{split}
b={}&(-495,-116,-684,108,601,-90,-173,-88,313,560,-397,461),\\
c={}&(-425,-281,-481,671,-255,30,475,-543,680,251,6,-128),\\
q={}&6263.
\end{split}
\]
In the preceding three orbits, the signed record is
\[
 \begin{aligned}
 U^{(1)}&=(2378,-452,-668,-322),\\
 U^{(2)}&=(1406,998,-204,-28),\\
 U^{(3)}&=(2479,-551,-371,-997).
 \end{aligned}
\]
Applying \(H_4/4\) to these three vectors reproduces
\eqref{vi:dep:K-testigo}. These identities verify the representation of
the incidences; their genesis remains that of the ledger
\eqref{vi:dep:incidencias}.

\subsection{\texorpdfstring{Periodic representation of \(K\) and normalisation of \(\alpha\)}{Periodic representation of K and normalisation of alpha}}
\label{vi:dep:alpha}

Fix \(B=1000\), the phase origin, and the twelve positions. Write
\[
 N_K=\sum_{j=1}^{12}K_jB^{12-j},\qquad
 \kappa=\frac{N_K}{B^{12}-1}.
\]
This is the exact periodic evaluation of the record. Multiplication
by \(B^{12}-1\) and successive Euclidean division recover its twelve
blocks, including any leading zeros. For
\(\rho K=(K_2,\ldots,K_{12},K_1)\), the same concatenation gives
\(\kappa(\rho K)=B\kappa(K)-K_1\). Thus reversibility
of the record preserves a specified phase origin.

\begin{proposition}[Minimal period of the rational representation]
\label{vi:dep:periodo-minimo-k}
The periodic prolongation of the record \eqref{vi:dep:K-testigo}
has minimal period \(12\) blocks in base \(1000\) and
\(36\) digits in base \(10\).
\end{proposition}
\begin{proof}
The twelve components of \(K\) are distinct. Any nonzero cyclic
translation smaller than twelve would identify two components if it were
a period. The minimal block period is therefore twelve.

Concatenating its three-digit blocks, including \(058\),
produces a decimal period of length \(36\). This expansion is
canonical: it is not eventually constant at the digit \(9\).
Let \(d\) be its minimal decimal period. Then \(d\mid36\),
because the shifts preserving a periodic word form
a cyclic subgroup. Grouping the digits in threes, the minimal
block period is the least \(k>0\) such that \(d\mid3k\), namely
\(d/\gcd(d,3)\). No proper divisor of \(36\)
satisfies \(d/\gcd(d,3)=12\). It follows that \(d=36\).
\end{proof}

This period belongs to the rational coordinate of the record, with the
phase origin fixed. The following composition also preserves the
regional coordinates and the normalisation quotients of its output.

Let \(P_j,E_j,\Phi_j\) be the fractional blocks of the three
regional characters already constructed, and extend \(K_j\) with period
twelve. The composition of records is
\begin{equation}
 Z_j=P_j+E_j-\Phi_j-K_j,
 \qquad
 A_j=Z_j+c_{j+1}-Bc_j.
 \label{vi:dep:normalizacion-alpha}
\end{equation}
The integer quotient \(c_j\) transports the continuation towards the observed
position. In a finite window, \(c_{N+1}\) is fixed first, and then
\[
 c_j=\left\lfloor\frac{Z_j+c_{j+1}}B\right\rfloor,
 \qquad A_j=Z_j+c_{j+1}-Bc_j.
\]
Euclidean division uniquely determines \(0\le A_j<B\).
Since \(-1998\le Z_j\le1998\), the set
\(\{-2,-1,0,1,2\}\) is stable for these quotients.

\begin{theorem}[Exact composition at unlimited depth]
\label{vi:dep:alpha-completa}
The complete evaluation of the normalisation, in the canonical decimal
convention, is
\begin{equation}
 \alpha_A=\pi+e-\varphi-4-\kappa.
 \label{vi:dep:alpha-A}
\end{equation}
There is a unique canonical expansion \((A_j)\), with blocks between zero
and \(B-1\), and a bounded sequence of integer quotients satisfying
\eqref{vi:dep:normalizacion-alpha} and \(c_1=0\).
\end{theorem}
\begin{proof}
The regional integer parts are \(3,2,1\). By absolute convergence,
\[
 y:=\sum_{j\ge1}Z_jB^{-j}
   =(\pi-3)+(e-2)-(\varphi-1)-\kappa.
\]
The first blocks are \(141,718,618,234\), so that
\(Z_1=7\). The absolute value of the tail sum is at most
\(1998\sum_{j\ge2}B^{-j}=0.002\); in particular, \(0<y<1\).
Take the canonical expansion \(y=\sum A_jB^{-j}\), choosing the
terminating expansion rather than the eventually constant \(B-1\) expansion if a
boundary occurs. Set
\[
 c_j=\sum_{r\ge0}(Z_{j+r}-A_{j+r})B^{-r-1}.
\]
Equality of the complete values shows that each \(c_j\) is the
negative of a finite sum of integers weighted by powers of
\(B\); it is therefore an integer, and \(c_1=0\). Separating the first term
gives \(Bc_j=Z_j-A_j+c_{j+1}\). The tail bounds make the
sequence bounded. More precisely, the tail of \(Z\) lies in \([-2,2]\)
and the canonical tail of \(A\) in \([0,1)\), so
\(-3<c_j\le2\). Being integral, \(c_j\) belongs to the stated set
\(\{-2,-1,0,1,2\}\).
Conversely, summing the recurrence through \(N\) gives
\[
 \sum_{j=1}^N A_jB^{-j}
 =\sum_{j=1}^N Z_jB^{-j}-c_1+c_{N+1}B^{-N}.
\]
Boundedness makes the final term vanish in the limit. Uniqueness of the
canonical expansion of \(y\) determines \(A\), and the tail formula
then determines every \(c_j\).
\end{proof}

The finite evaluation \(\kappa_{36}=N_K/B^{12}\) differs from
\(\kappa\) by \(B^{-12}\kappa\). Moreover, imposing
\(c_{13}=0\) in a twelve-block window does not replace the quotient
of the complete continuation. In the expressed record, the continuation
gives \(c_{13}=-1\): the twelfth block is \(662\), whereas the
truncated window with zero terminal quotient gives \(663\). The complete
evaluation begins
\[
 \alpha_A=0.007297352569283800997285105472380662999\ldots.
\]
The period of \(K\) does not impose a period on \((A_j)\), since
\eqref{vi:dep:normalizacion-alpha} also involves the three
regional records and the integer quotients.

In the following formulas, \(\alpha\) denotes this complete positional
output \(\alpha_A\). The analytic vacancy route has its own
operator; a polynomial truncated at order nine and the root of that polynomial
are not identified here with the complete operator. The mass
compositions that follow do not require that approximation to be promoted to an
equality between the two infinite routes.

\subsection{Three-dimensional construction of the electronic prefactor}
\label{vi:dep:prefactor}

Let \(\mathcal A_3=\{1,\ldots,9\}^3\), with uniform measure, and
\(\rho_9(n)=1+((n-1)\bmod9)\). On this domain act
\[
 \Sigma(x)=\rho_9(x_1+x_2+x_3),\qquad
 \Pi(x)=\rho_9(x_1x_2x_3).
\]
On \(\ell^2(\mathcal A_3)\), let \(P\) and \(Q\) be the conditional
averages over the fibres of \(\Sigma\) and \(\Pi\),
respectively. Each is idempotent and self-adjoint: within each
fibre it replaces a vector by its constant component and preserves
that component orthogonally. They are therefore orthogonal projectors.

The additive fibres have cardinality \(81\). The multiplicative fibres have
cardinalities
\[
m=(36,36,108,36,36,108,36,36,297).
\]
The joint incidence
\(N_{ab}=\#\{x:\Sigma(x)=a,\Pi(x)=b\}\) is \(9N_0\), where
\begin{equation}
 N_0=\begin{pmatrix}
0&1&1&0&1&1&0&1&4\\1&0&1&1&0&1&1&0&4\\
1&0&2&0&1&2&0&0&3\\0&1&1&0&1&1&0&1&4\\
1&0&1&1&0&1&1&0&4\\0&0&2&1&0&2&0&1&3\\
0&1&1&0&1&1&0&1&4\\1&0&1&1&0&1&1&0&4\\
0&1&2&0&0&2&1&0&3
\end{pmatrix}.
 \label{vi:dep:incidencia-cubica}
\end{equation}
This matrix is obtained by counting the \(729\) triples. It is also
reproduced without rounding by
\(N_{ab}=\sum_{i,j,k=1}^9
\mathbf1_{\{a\}}(\rho_9(i+j+k))
\mathbf1_{\{b\}}(\rho_9(ijk))\).
In the normalised bases of functions constant on each fibre, the
intertwining has matrix \(C_{ab}=N_{ab}/\sqrt{81m_b}\).

\begin{proposition}[Norm of the noncommutativity of the two sheets]
\label{vi:dep:norma-PQ}
The preceding projectors satisfy
\begin{equation}
 \|[P,Q]\|=\frac{\sqrt3}{4}.
 \label{vi:dep:prefactor-valor}
\end{equation}
\end{proposition}
\begin{proof}
On the space of the nine additive fibres, \(PQP\) has matrix
\(M=CC^{\mathsf T}\). Multiplication of the matrices in
\eqref{vi:dep:incidencia-cubica} gives eigenvalue one on the constant
vector, \(1/4\) on \((1,-1,0,1,-1,0,1,-1,0)\), and
\(7/132\) on \((1,1,-2,1,1,-2,1,1,-2)\).
On coordinates \(\{3,6,9\}\) with zero sum it acts by
\(1/18\), with multiplicity two. On each of the sets
\(\{1,4,7\}\) and \(\{2,5,8\}\), the zero-sum vectors
belong to the kernel; they give four dimensions. These orthogonal
subspaces add up to the nine dimensions.

For a unit eigenvector \(u\in\operatorname{ran}P\) of
\(PQP\) with eigenvalue \(0<\lambda<1\), the vector
\(v=(Qu-\lambda u)/\sqrt{\lambda(1-\lambda)}\) is a unit vector
orthogonal to \(\operatorname{ran}P\). In the basis \((u,v)\),
\[
 P=\begin{pmatrix}1&0\\0&0\end{pmatrix},\quad
 Q=\begin{pmatrix}\lambda&\sqrt{\lambda(1-\lambda)}\\
 \sqrt{\lambda(1-\lambda)}&1-\lambda\end{pmatrix}.
\]
The commutator has norm \(\sqrt{\lambda(1-\lambda)}\).
The blocks corresponding to \(\lambda=0,1\) commute. Among the
calculated eigenvalues, the maximum is attained at \(\lambda=1/4\),
since \(\lambda(1-\lambda)\) is increasing on \([0,1/2]\).
Its square root is \(\sqrt3/4\).
\end{proof}

In the extremal block, \(S=2P-I\) and
\(J=4[P,Q]/\sqrt3\) satisfy
\[
 S^2=I,\qquad J^2=-I,\qquad SJ=-JS.
\]
The four matrices \(I,S,J,SJ\) are linearly independent,
so they realise \(\mathrm{Cl}_{1,1}\simeq M_2(\RR)\).
The elements \(n_\pm=(S\pm J)/2\) satisfy
\(n_\pm^2=0\) and \(n_+n_-+n_-n_+=I\), by direct expansion.
The prefactor of the electronic equation thus arises from two evaluations
of APP in dimension three and their commutator, before any reference
mass.

\subsection{Boundary reading and basal electronic composition}
\label{vi:dep:basal}

The incidence \eqref{vi:dep:incidencia-hojas} has positive eigenvector
\(r=(1,\varphi)^{\mathsf T}\). The normalisation
\(p_{ij}=F_{ij}r_j/(\varphi r_i)\) produces
\[
 P_{\rm av}=\begin{pmatrix}0&1\\\varphi^{-2}&\varphi^{-1}\end{pmatrix},
 \qquad \mu=\frac{(1,\varphi^2)}{1+\varphi^2}.
\]
The identity \(\varphi^{-2}+\varphi^{-1}=1\) proves that the rows
sum to one. Since \(F\) is symmetric, weights proportional to
\(r_i^2\) satisfy detailed balance; hence \(\mu P_{\rm av}=\mu\).
The incidence is irreducible and has a self-loop, so
this stationary distribution is unique. The chronological frequency
\((1/3,2/3)\) counts the three instants of the calendar; \(\mu\)
weights its admissible trajectories with memory one.

The operator preserving the sheets, fixing the volumetric normalisation
at one, and marking the areal sheet by the closure character is
\(D_\pi=\operatorname{diag}(\pi,1)\). By induction,
\[
 F^n=\begin{pmatrix}F_{n-1}&F_n\\F_n&F_{n+1}\end{pmatrix}.
\]
The ratio of marked diagonal returns is
\[
 \Theta_n=\pi\frac{F_{n-1}}{F_{n+1}}
 \longrightarrow\frac{\pi}{\varphi^2}.
\]
Convergence follows from the dominant eigenvalue \(\varphi\), or from
the consecutive ratios already bounded. Under interchange of its arguments, the reader
\(\mathcal R(x,y)=x/y^2\)
also expresses
\begin{equation}
 \rho_A=\frac{\pi}{\varphi^2},\qquad
 \rho_E=\frac{\varphi}{\pi^2},\qquad
 \varphi^3\rho_A^2\rho_E=1.
 \label{vi:dep:dos-orientaciones}
\end{equation}
The last equality is checked by multiplying the fractions.
It relates two orientations of the same reader and does not select the
remaining coefficients of the electronic equation.

The typed central return has length \(27\), with support
\(\ZZ/9\ZZ\times\FF_3\). In the fixed regional gauge, the five microfibre
positions are inserted as
\((4,0),(5,0),(6,0),(6,1),(6,2)\).
The complement has \(27-5=22\) positions. Its incidence with the
three tritic phase states has \(5\cdot3=15\) pairs.
This yields the integers \(-22\) and \(15\) used by the
electronic reader; orientation fixes the sign of the first term.
The cubic composition uses three copies of the closure coordinate,
so its evaluation is \(\alpha^3\). The support, regional
insertion, and this composition rule are distinct operational data;
an isolated count does not replace their specification.

The global defect and its nonadic return are expressed through
\begin{equation}
 d_4=\frac{\pi-e\log\pi}{270},\qquad
 \Delta_4=d_4\left(1+\frac{\pi}{729}\right).
 \label{vi:dep:delta4}
\end{equation}
The denominator \(270\) belongs to the corona without its closure
unit; \(729=9^3\) belongs to the cubic support. The function
\(f(x)=x-e\log x\) attains its minimum zero at \(x=e\), since
\(f'(x)=1-e/x\) changes sign there. As \(\pi\ne e\),
we obtain \(d_4>0\) and \(\Delta_4>0\).

\begin{definition}[Basal electronic coordinate]
\label{vi:dep:electron-basal}
The composition of the preceding readers is
\begin{equation}
 \varepsilon_e^{(0)}=
 \frac{\sqrt3}{4}
 \left[\exp\left(\frac{\varphi}{\pi^2}\right)-22\alpha^3\right]
 (1+15\Delta_4).
 \label{vi:dep:formula-basal}
\end{equation}
\end{definition}
The exponential acts on orientation \(\rho_E\) of
\eqref{vi:dep:dos-orientaciones}; the vacancy reading corrects that
component, and the final factor preserves the return of the global defect.
For \(0<\alpha<0.01\), all factors are positive:
\(\exp(\varphi/\pi^2)>1\), \(22\alpha^3<22\cdot10^{-6}<1\),
and \(\Delta_4>0\). This is a dimensionless coordinate. Its presence
in a mass chart is defined after construction of the scale.

\subsection{Determinantal scale and the reduced Planck constant}
\label{vi:dep:planck}

The local block \eqref{vi:dep:hadamard} has determinant \(-16\).
Its Smith normal form is \(\operatorname{diag}(1,2,2,4)\).
Indeed, the greatest common divisor of the entries is one; the minors of
order two are even, and one has absolute value two. By the adjugate identity,
the minors of order three have absolute value four,
and the determinant has absolute value sixteen. The determinantal
divisors are \(1,2,4,16\), whose successive ratios give the
four Smith factors. Consequently,
\begin{equation}
 G_H=\operatorname{coker}H_4
 \simeq\ZZ/2\ZZ\oplus\ZZ/2\ZZ\oplus\ZZ/4\ZZ,
 \qquad |G_H|=16.
 \label{vi:dep:conucleo}
\end{equation}

The local normed reader multiplies the same coordinate \(\alpha\)
once for each class of this cokernel and applies the self-scaling
coordinate once. Its definition produces
\begin{equation}
 \Sigma_H=\varphi\prod_{g\in G_H}\alpha
         =\varphi\alpha^{16}.
 \label{vi:dep:sigma-H}
\end{equation}
The order of the cokernel determines the exponent of this reader. The Smith
normal form and the choice of normed reader serve distinct
functions: the former does not turn an arbitrary functional of the block into
expression \eqref{vi:dep:sigma-H}.

\begin{proposition}[Decade of action]
\label{vi:dep:decada}
The output \eqref{vi:dep:sigma-H} satisfies
\(10^{-34}<\Sigma_H<10^{-33}\); its decimal valuation is \(-34\).
\end{proposition}
\begin{proof}
The regional and positional evaluations give the intervals
\(1.618<\varphi<1.619\) and \(0.007297<\alpha<0.007298\).
By positivity,
\[
 \frac{1618\,7297^{16}}{10^{99}}
 <\Sigma_H<
 \frac{1619\,7298^{16}}{10^{99}}.
\]
The integer inequalities
\(1618\,7297^{16}>10^{65}\) and
\(1619\,7298^{16}<10^{66}\) prove the result without using an
external action constant.
\end{proof}

The decade and the mantissa are representations of different stages.
Before constructing the latter, the regional reader selects the persistent
return of the closure microfibre. The projection
\(p_{\rm ret}(u_1,\ldots,u_6)=(u_4,u_5,u_6)\), applied to the
five emissions of \ref{vi:dep:regiones}, produces the block
\(b_\pi=(169,272,272)\) four times and
\(b_\pi^-=(418,674,521)\) once. These multiplicities count
regional rows, not the seed multiplicities \(432,144\).

\begin{proposition}[Equivariant selector of the regional return]
\label{vi:dep:selector169}
Among rules preserving permutations of the five regions,
selecting the terminal block of maximal multiplicity, equivariant
under permutations of its three positions, and retaining a position fixed
by the stabiliser of the block, the selected coordinate is unique
and equals \(\rho_\pi=169\).
\end{proposition}
\begin{proof}
Projection of the five rows shows that the block of maximal
multiplicity is unique and is \(b_\pi\). Its stabiliser in
\(\mathfrak S_3\) interchanges the two positions of value \(272\)
and fixes only the position of value \(169\). Retaining that position
is compatible with any simultaneous permutation of coordinates.
Equivalently, the value is expressed without choosing an initial position
as
\(\operatorname{sing}(b_\pi)=\sum_j(b_\pi)_j
-2\operatorname{moda}(b_\pi)=169\).
\end{proof}

Grading by length expresses a word of length \(n\)
through the character \(\chi_z([w])=z^{|w|}\). Concatenation
adds lengths, so
\(\chi_z([uv])=\chi_z([u])\chi_z([v])\). Since the regional
emissions belong to \(\mathcal U_6\), the return impulse
appears as \(\rho_\pi z^6\). Degree six corresponds to the
length of the emission, not the length of a closed walk
traversing several anchors of the graph.

The prolongation weight is obtained on the determinant line of
bilayer transport. For a real involution \(R\) in dimension two,
with \(\operatorname{tr}R=0\), let
\[
 x_A=\frac{50\pi}{9}\alpha,\qquad
 \mathcal T(x_A,y)=\exp(-x_AI_2-yR).
\]
The oriented coordinate \(y\) is still formal. In an eigenbasis
of \(R\), the two entries are \(e^{-x_A-y}\) and
\(e^{-x_A+y}\). Therefore,
\[
 \bigwedge^2\mathcal T=\det\mathcal T
 =e^{-2x_A}=D_A.
\]
Dependence on \(y\) cancels. Multiplicativity under
concatenation determines the prolongation character \(k\mapsto D_A^k\):
any character with elementary value \(D_A\) must have that value
at \(k=1+\cdots+1\). This construction precedes evaluation
of the conjugate angular coordinate.

The fifth-order regional reader uses eigenvalues \(9,5\)
of the central block \(\bigl(\begin{smallmatrix}7&2\\2&7\end{smallmatrix}\bigr)\),
the orbit of four positions, and the calendar of twelve windows.
Its graded rule fixes
\begin{align}
 H_5={}&\frac12\sqrt{\frac{1000\alpha}{\varphi}}-\alpha
       +\frac9{16}\alpha^2-\frac59\alpha^3
       +\frac7{48}\alpha^4-\frac1{54}\alpha^5,\notag\\
 D_A={}&\exp\left(-\frac{100\pi\alpha}{9}\right),\notag\\
 \mathcal C_\pi={}&169\alpha^6+
                  \frac{D_A\alpha^7}{1-D_A\alpha},\notag\\
 \eta_{\rm ret}={}&H_5-\frac{90}{\pi}\mathcal C_\pi.
 \label{vi:dep:mantisas}
\end{align}
In this rule, the ratios are
\(9/16=9/4^2\), \(5/9\), \(7/48=(\operatorname{tr}B_c/2)/(4\cdot12)\),
and \(1/54=1944/104976\). Their assignment to the degrees and signs
forms part of the graded reader. The integer \(169\) is the representation
of the selector proved in Proposition \ref{vi:dep:selector169}.
Its regional datum precedes the evaluation of \(\alpha\); the
analytic reader subsequently composes it with the closure coordinate.

The tail of \(\mathcal C_\pi\) has a renewal construction.
The impulse \(\rho_\pi z^6\) and the strict continuation belong
to different summands. After recording the impulse once,
the continuation is normalised by the unit of the determinant line;
each new step multiplies it by \(D_A\). The orientation of the
regional gauge places this return in the compensating sector
subtracted from the fifth-degree character.
With formal parameter \(z\), the rule
\begin{equation}
 R(z)=D_Az^7+D_AzR(z)
 \label{vi:dep:renovacion}
\end{equation}
uniquely determines
\(R(z)=D_Az^7/(1-D_Az)\). This can be proved by comparing
coefficients or by inverting the formal series \(1-D_Az\).
Normalisation of the first return, in addition to the recurrent ratio,
is essential to this uniqueness. Since \(0<D_A<1\) and \(\alpha<1\),
evaluation at \(z=\alpha\) converges absolutely.

\begin{proposition}[Positivity and the action return ratio]
\label{vi:dep:positividad-accion}
On the generated branch, \(0<\eta_{\rm ret}<H_5\). The ratio
\begin{equation}
 R_{\rm act}=\frac{H_5}{\eta_{\rm ret}}
 \label{vi:dep:Ract}
\end{equation}
is positive, greater than one, and dimensionless.
\end{proposition}
\begin{proof}
The bounds \(0.007<\alpha<0.008\) and \(\varphi<1.619\) give
\(\tfrac12\sqrt{1000\alpha/\varphi}>1\). The sum of the absolute
values of the remaining negative terms of \(H_5\) is less than
\(0.008001\), so \(H_5>0.99\).
Since \(D_A<1\), \(\pi>3\), and \(\alpha<0.008\),
\[
 0<\frac{90}{\pi}\mathcal C_\pi
 <30\left(169(0.008)^6+
                  \frac{(0.008)^7}{1-0.008}\right)<10^{-6}.
\]
It follows that \(0<\eta_{\rm ret}<H_5\). The ratio of two action
coordinates in the same basis is dimensionless and does not change when
that basis is rescaled.
\end{proof}

Let \(\mathcal U_S\) be a basis of the action line. The two
sections, before and after the return, are
\begin{equation}
 \hbar_{\rm pre}=H_5\,10^{-34}\mathcal U_S,
 \qquad
 \hbar_{\rm ret}=\eta_{\rm ret}\,10^{-34}\mathcal U_S.
 \label{vi:dep:hbar}
\end{equation}
Their physical reading is the proposed construction of the reduced Planck
constant, with the return section made explicit. The unit
\(\mathcal U_S\mapsto\mathrm{J\,s}\) is a subsequent metrological
realisation. The coefficient of \(\Sigma_H\) within its decade is not
\(H_5\): the determinant expresses the scale, and the regional reader
expresses the coordinate occupying that scale. For each section,
\(h=2\pi\hbar\) expresses the action of a complete circuit.

\subsection{Angular coordinates and the clock of the absolute realisation}
\label{vi:dep:reloj}

The angular pair uses the dimensionless coordinates already generated:
\begin{equation}
 A_{\deg}=1000\alpha,\qquad
 C^*_{\deg}=2(\eta_{\rm ret}+\alpha),\qquad
 x=\frac{\pi}{180}A_{\deg},\quad
 y=\frac{\pi}{180}C^*_{\deg}.
 \label{vi:dep:angulos}
\end{equation}
The first operation shifts the base-thousand expansion by one block.
The second is the bidirectional representation of the regional coordinate:
equivalently,
\(y=\pi(H_5+\alpha)/90-\mathcal C_\pi\).
This equality follows by substituting \eqref{vi:dep:mantisas}; it
introduces no additional rule for selecting the counter-angle.
The chart of one circuit has \(360\) angular units or \(2\pi\)
radians. Thus, for a lifted coordinate,
\[
 \hbar\theta_{\rm rad}
 =h\frac{\theta_{\deg}}{360}.
\]
The factor \(10^{-34}\) remains on the action line and is not
inserted into the angular coordinates.

The positive channels are
\[
 q_+=e^{-x-y},\qquad q_-=e^{-x+y}.
\]
The pair is recovered through
\(x=-\tfrac12\log(q_+q_-)\),
\(y=\tfrac12\log(q_-/q_+)\). Channel interchange preserves
\(x\) and reverses the sign of \(y\). The operator with oriented
projectors \(P_\pm\),
\(T=q_+P_++q_-P_-\), preserves those labels when publishing its two
eigenvalues.

For an internal time unit \(t_0>0\), the closure of \(108\)
steps determines
\begin{equation}
 \omega_0=\frac{2\pi}{108t_0},\qquad
 E_{\rm clk}=\hbar_{\rm ret}\omega_0,
 \qquad m_{\rm clk}=\frac{E_{\rm clk}}{c_{\rm int}^{2}}.
 \label{vi:dep:reloj-absoluto}
\end{equation}
These formulas produce dimensions of frequency, energy, and mass.
Velocity and length are obtained through the response of the
same channels, as specified by the following identities. Expressing them
in seconds, joules, or mass units requires the corresponding
charts; those charts do not alter the channels or the electronic
section normalised relative to the clock.

\subsubsection{Identity between constitutive velocity and clock velocity}

In the domain \(x>|y|\), the two channels satisfy \(0<q_\pm<1\).
On the two-sheet realisation, with rank-one projectors
\(P_+,P_-\), the temporal response is defined by
\[
 \mathsf V=(I-T^{90})(I-T^{120})^{-1}
          =r_+P_++r_-P_-,
 \qquad r_\pm=\frac{1-q_\pm^{90}}{1-q_\pm^{120}}>0.
\]
The inverse exists because both eigenvalues of \(T^{120}\) are less
than one. The symmetric reading of the response is
\[
 \mathcal E=
 \log[(1-q_+^{90})(1-q_-^{90})]
 -\log[(1-q_+^{120})(1-q_-^{120})]
 =\log(r_+r_-).
\]
Let \(u_S,u_C,u_T\) be positive bases of action, velocity, and
time, and \(u_L=u_Cu_T\) the length basis. The kinematic
and constitutive readings are
\[
 c_{\rm int}=e^{-\mathcal E}u_C,\qquad
 c_{\rm vac}=(r_+r_-)^{-1}u_C.
\]

\begin{proposition}[Common realisation of propagation and absolute mass]
\label{vi:dep:velocidad-comun}
In the same dimensional chart,
\[
 c_{\rm int}=c_{\rm vac}=\widehat c\,u_C,
 \qquad \widehat c=(\det\mathsf V)^{-1}=(r_+r_-)^{-1}.
\]
For \(t_0=\tau_0u_T\), with \(\tau_0>0\), the elementary length is
\[
 \ell_0=c_{\rm int}t_0=\tau_0\widehat c\,u_L.
\]
The chronological mass of \eqref{vi:dep:reloj-absoluto} therefore
retains the form
\[
 m_{\rm clk}=
 \frac{2\pi\,\eta_{\rm ret}\,10^{-34}}
      {108\,\tau_0\,\widehat c^{\,2}}\,
 \frac{\mathcal U_S}{u_Tu_C^2}.
\]
\end{proposition}
\begin{proof}
The identity \(\mathcal E=\log(r_+r_-)\) gives
\(e^{-\mathcal E}=(r_+r_-)^{-1}\); the rank-one property of each sheet gives
\(\det\mathsf V=r_+r_-\). Sheet interchange permutes the
factors and preserves their product. Length is obtained by multiplying
the same velocity by \(t_0\). Substituting these expressions and
\eqref{vi:dep:hbar} into \eqref{vi:dep:reloj-absoluto} produces the
mass formula. The factor \(\mathcal U_S/(u_Tu_C^2)\) has
the dimension of mass; it does not enter the selection of the channels.
\end{proof}

For a homogeneous route of \(n>0\) transitions, the additive readings
\(T(\gamma)=nt_0\) and \(L(\gamma)=n\ell_0\) give
\(L(\gamma)/T(\gamma)=c_{\rm vac}\). If the channel varies from
edge to edge, their lengths are added; the ratio is then the mean
velocity of those edges, not necessarily a common local velocity.

Equality of sections persists under changes of chart. If
\(u_C'=du_C\), \(u_T'=\eta u_T\), with \(d,\eta>0\), then
\[
 \widehat c'=\widehat c/d,\qquad
 \tau_0'=\tau_0/\eta,\qquad u_L'=d\eta u_L,
 \qquad
 \tau_0'\widehat c'u_L'=\ell_0.
\]
In particular, choosing the adapted basis \(u_C'=c_{\rm vac}\)
makes \(\widehat c'=1\), but the reader coordinate in the
internal chart remains \((r_+r_-)^{-1}\).

\subsection{Two readers of the electronic trajectory and beta evaluation}
\label{vi:dep:beta}

The central trajectory starts at \((r,c)=(5,5)\), with
orientation \((h,v)=(1,1)\). TRIT repeats the cycle \((+1,-1,0)\).
In the first regime, \(r+c\) is read and the advance
\(c\mapsto1+((c-1+h)\bmod9)\) follows; in the second, \(rc\) is read and
the advance \(r\mapsto1+((r-1+v)\bmod9)\) follows; in the zero
regime, \(9\) is read without displacement. The class modulo
three of the digital root of the reading is expressed. After steps \(27,54,81\),
\(h,v\) are simultaneously reversed. The representation of \(108\)
steps preserves the complete trajectory and gives
\begin{equation}
 a=100000220,\quad b=120200000,\quad
 \gamma_e=(a^3b^3)^2,
 \quad\tau_e=(+,+,+,-,-,-,+,+,+,-,-,-).
 \label{vi:dep:trayectoria-electron}
\end{equation}
The identity is checked by traversing the nine instants of each
window; at the third block, reversal changes \(a\) to \(b\),
and the second traversal of \(54\) steps repeats the representation.
The state also preserves its memory coordinates.

The first reader measures cyclic disagreements of the ternary representation:
\[
 D_k=\#\{t\in\ZZ/108\ZZ:\gamma_e(t)\ne\gamma_e(t+k)\}.
\]
Since the expressed word has period \(54\), every disagreement
appears together with its translate by \(54\); all \(D_k\) are even.
In the half \(a^3b^3\), the disagreement count for
\(k=1,3,4,27,36\) gives respectively \(27,28,34,24,16\).
Doubling it yields
\[
 (D_1,D_3,D_4,D_{27},D_{36})=(54,56,68,48,32).
\]
Hence
\begin{equation}
 K_D=\left(D_{27}+D_{36},D_1,\frac{D_4-D_3}{2}\right)
     =(80,54,6).
 \label{vi:dep:KD}
\end{equation}

The second reader acts on the orientation of the twelve windows.
Define, with cyclic indices,
\[
 C_k=\sum_{r=0}^{11}\tau_r\tau_{r+k},\quad
 A_\ell=\sum_{r=0}^{11}\left(\sum_{j=0}^{\ell-1}\tau_{r+j}\right)^2,
 \quad P_6=\sum_{r=0}^{5}\tau_r\tau_{r+6}.
\]
The oriented primitive integer combination cancelling the diagonal part
of \(A_3,A_4\) is \(\Omega=4A_3-3A_4\): its coefficients
satisfy \(3u+4v=0\), \(\gcd(u,v)=1\), and \(u>0\) is fixed.
Expanding the square proves
\[
 A_\ell=\ell C_0+
     2\sum_{k=1}^{\ell-1}(\ell-k)C_k,
 \qquad\Omega=-2C_1-4C_2-6C_3.
\]
For \(\tau_e\), we obtain
\(C_0=12,C_1=4,C_2=-4,C_3=-12\),
\(A_3=44,A_4=32,P_6=6\). Therefore,
\begin{equation}
 K_\Omega=(\Omega,54,P_6)=(80,54,6)=K_D.
 \label{vi:dep:KOmega}
\end{equation}
The second component preserves the half-return of \(108\) steps;
it is not introduced by diagonal cancellation. Equality of the two
readers has been proved on the specified central trajectory;
it is not extended by definition to all collections of the atlas.

\begin{definition}[Electronic coordinate after the return]
The complete electronic composition is
\begin{equation}
 \varepsilon_e^\beta
  =\varepsilon_e^{(0)}
      R_{\rm act}^{80-54\alpha+6\Delta_4}.
 \label{vi:dep:electron-beta}
\end{equation}
\end{definition}
Positivity of \eqref{vi:dep:formula-basal} and
\eqref{vi:dep:Ract} allows this power to be defined through the
real logarithm. The integer triple arises equally from the two
central readers just calculated. Relative evaluation
preserves the reduced Planck constant through the ratio of its
sections, and absolute realisation additionally uses
\eqref{vi:dep:reloj-absoluto} and the stated dimensional chart.

In the inherited tables, \(\mathfrak e_0\) denotes the
basal coordinate \(\varepsilon_e^{(0)}\) in the chart
specified there. The section \(\hbar_{\rm int}\) used for a
survival evaluation must be stated in the same chart:
for the post-return realisation of the present development,
it is \(\hbar_{\rm ret}\). Historical rows retain the
convention stated in their source table and are not recalculated through a
silent change of this section.

If \(\mathcal U_E\) is the energy basis of the electronic
evaluation, \(E_e=\varepsilon_e^\beta\mathcal U_E\). When
the chronological quantum is used as the basis, the coefficient is
\[
 \widehat\varepsilon_e
   =\varepsilon_e^\beta\frac{\mathcal U_E}{E_{\rm clk}},
 \qquad E_e=\widehat\varepsilon_e E_{\rm clk}.
\]
The equality preserves the object and transforms its coordinate. Hence
a number presented in an MeV chart cannot be multiplied
without conversion by an arbitrary different energy quantum. The electron
retains its structural construction when expressed relative to a
chronological scale: it is not replaced by the mass of that quantum.

Metrological comparison applies to the result of this composition,
with the version of \(\alpha\), the action section, the
\(\Delta_4\) reader, and the unit chart stated. In particular,
\(d_4\) and \(\Delta_4\) are not interchangeable: replacing them
simultaneously in the basal factor and in the exponent changes the
result by the positive ratio
\[
 \frac{1+15\Delta_4}{1+15d_4}
 R_{\rm act}^{6(\Delta_4-d_4)}>1.
\]
The preserved formulas allow that difference to be located without
attributing it to an external mass or selecting the trajectory anew.
